%%%PAGE 214 rot=0 legibility=good summary=伏安特性曲线(电阻丝、半导体、标准电阻的U-I与I-U图像)；伏安特性曲线与电源U-I图像的交点问题(单灯、双灯串联、双灯并联)
%%%BLOCK id=214-01 ch=10 sec=伏安特性曲线·电阻与电阻丝 kind=knowledge alt=none cont=none y=0.07-0.30
\textbf{伏安特性曲线}

%FIG p214_a 0.09 0.12 0.28 0.22
\notefig[0.25]{figures/p214_a.png}

%FIG p214_b 0.29 0.12 0.49 0.22
\notefig[0.25]{figures/p214_b.png}

电阻/电阻丝

偏$U$

$I\uparrow$ \unclear{$P\uparrow$} $T\uparrow$ $\rho\uparrow$ $R\uparrow$

\noindent 割线斜率 $k=\dfrac{U}{I}=R$ \qquad\qquad $k=\dfrac{I}{U}=\dfrac{1}{R}$\\[4pt]
矩形面积 $=U\cdot I=P_{\text{瞬时}}$ \qquad\quad $S=U\cdot I=P_{\text{瞬时}}$。
%%%END
%%%BLOCK id=214-02 ch=10 sec=伏安特性曲线·半导体 kind=knowledge alt=none cont=none y=0.31-0.46
半导体 \quad 温度$\uparrow$ $R\downarrow$。

%FIG p214_c 0.08 0.35 0.485 0.455
\notefig[0.45]{figures/p214_c.png}
%%%END
%%%BLOCK id=214-03 ch=10 sec=伏安特性曲线·标准电阻 kind=knowledge alt=none cont=none y=0.49-0.62
标准电阻 \quad $\rho$不变

%FIG p214_e 0.08 0.515 0.28 0.62
\notefig[0.25]{figures/p214_e.png}

%FIG p214_f 0.285 0.515 0.49 0.62
\notefig[0.25]{figures/p214_f.png}
%%%END
%%%BLOCK id=214-04 ch=10 sec=伏安特性曲线与电源U-I图像交点 kind=method alt=none cont=none y=0.65-0.83
伏安特性曲线与电源 $U$-$I$ 图像的交点问题（可变电阻求功率）。\quad 注意等效电源内阻/电源电动势

%FIG p214_g 0.085 0.72 0.56 0.815
\notefig[0.55]{figures/p214_g.png}

\[ E=U+Ir \]
\[ U=-Ir+E \]
%%%END
%%%BLOCK id=214-05 ch=10 sec=伏安特性曲线与电源U-I图像交点 kind=method alt=none cont=none y=0.83-1.00
%FIG p214_i 0.07 0.832 0.235 0.90
\notefig[0.25]{figures/p214_i.png}

%FIG p214_j 0.24 0.835 0.425 0.955
\notefig[0.25]{figures/p214_j.png}

\[ E=2U+Ir \]
\[ U=-\frac{Ir}{2}+\frac{E}{2} \]
%%%END
%%%BLOCK id=214-06 ch=10 sec=伏安特性曲线与电源U-I图像交点 kind=method alt=none cont=none y=0.80-1.00
%FIG p214_k 0.47 0.795 0.87 0.935
\notefig[0.4]{figures/p214_k.png}

\[ E=U+2I\cdot r \]
\[ U=-2r\cdot I+E \]
%%%END
%%%PAGEEND 214
